\section{Discussion}
\label{sec:discussion}

\subsection{Why an alignment objective should matter more as sparsity rises}

The design hypothesis behind BaCP is not that a contrastive term is a better loss
than cross-entropy in general. It is that the two losses differ in \emph{how much
they constrain}, and that the difference is worth increasingly more as the
surviving parameter budget shrinks.

Cross-entropy on a $C$-way task supplies at most $\log_2 C$ bits per example, and
in practice far fewer once the model is confident: on CIFAR-10 a correctly and
confidently classified image contributes almost no gradient. Everything the
network computes before the final logits is unconstrained except through that
scalar. At 90\% sparsity this is adequate, because the surviving weights still
span a subspace close enough to the dense one that the label is enough to
re-tune the decision boundary within it. At 99.9\% it is not: the mask has removed
so much of the support that the sparse network is not re-tuning a representation,
it is choosing a new one, and a per-example scalar is a very weak prior over that
choice.

Equation~\eqref{eq:functional} supplies a different quantity. Its gradient with
respect to an anchor embedding is
$\tau^{-1}\big[\sum_k p_{ik}\hat z_k - |P(i)|^{-1}\sum_{j\in P(i)} \hat z_j\big]$,
the difference between the softmax-weighted centroid of all candidates and the
centroid of the positives: a $D$-dimensional target direction per example rather
than a scalar, computed against a frozen dense teacher that has not been damaged.
The teacher receives no gradient, so the term is a distillation signal --- it
transfers geometry the sparse network can no longer derive on its own. That is why
the effect should scale with damage, and why we extended the grid to 0.999
rather than stopping where the antecedent stopped.

The three teachers correspond to three failure modes, which is CAP's insight
rather than ours. PrC guards against losing the general features the ImageNet
initialisation supplied; FiC guards against losing the task-specific structure the
dense fine-tune built; SnC guards against trajectory drift, in which each mask
update pushes the network somewhere the previous updates did not anticipate. The
last of these is specific to \emph{gradual} pruning, and is the one that should
matter most at extreme sparsity, where the final increments of the cubic ramp
remove the largest fraction of what is still alive.

We flag one thing this account does not explain. The ablation in
Section~\ref{sec:results:ablation} found that the formulation with the larger,
looser similarity matrix and the trainable head outperformed the tighter
CAP-faithful one, which the argument above does not predict --- if anything it
predicts the opposite, since the cleaner formulation constrains the backbone more
directly. We report the ordering because we measured it, and we do not have a
mechanism for it. The most we can say is that the square-matrix formulation adds
within-student terms that the rectangular one lacks, which makes it a strictly
larger objective, and that the two effects are confounded in a two-variable
ablation. Separating them is future work.

\subsection{Limitations}
\label{sec:discussion:limitations}

\paragraph{One dataset, one modality.}
Every result in this paper is CIFAR-10. The claim that the objective transfers
from language models to vision backbones therefore rests on a single vision
dataset at a single resolution with a single label granularity. CIFAR-10 is also
a regime where dense accuracy is high and the headroom between arms is small,
which makes it a demanding setting for detecting a difference but a weak setting
for establishing generality. Nothing here shows the effect survives to CIFAR-100,
to ImageNet-scale data, or to the transformer backbones that our framework also
supports but that we did not include in this grid. A reader should treat the
result as evidence about this objective on this dataset family, and the
vision-transfer claim as demonstrated on one point rather than characterised.

\paragraph{The trainable projection head, and the claim it does not support.}
The objective is computed on $z = \mathrm{normalize}(g(f_\theta(x)))$, and in the
configuration we ran, $g$ is a trainable linear head of roughly $262{,}000$
parameters that is never pruned. Because the teacher branch is frozen, the
optimiser can lower $\mathcal{L}_{\mathrm{PrC}}$, $\mathcal{L}_{\mathrm{FiC}}$ and
$\mathcal{L}_{\mathrm{SnC}}$ by adjusting $g$ alone, leaving $\theta$ unchanged:
the minimisation admits solutions that never touch the backbone. Ordinary
contrastive learning is immune to this because both branches share the encoder, so
the only route to a lower loss is a better encoder; freezing one branch inserts a
free adapter between the optimised weights and the loss.

The consequence is a bound on the \emph{claim}, not on the \emph{measurement}.
The accuracies we report are measured after a full fine-tune of the masked
network with the projection head discarded, so they are legitimate results for a
deployed sparse model, and the matched-budget design means the difference between
arms is still attributable to the objective. What we cannot claim on the strength
of the loss alone is that BaCP \emph{preserves representations} --- that would
require the loss to be reducible only through $\theta$, and here it is not. The
defensible statement is the measured one: this objective produces better sparse
models under a matched budget. We ran the variant that would license the stronger
claim --- one frozen projection head shared by student and teachers, under which
the head carries no gradient and the backbone is the only route --- and it was worse
at every measured sparsity (Table~\ref{tab:ablation}). We report that honestly:
the configuration with the clean interpretation is not the configuration with the
better numbers, and we chose the numbers.

\paragraph{Compute budget below the published comparators.}
Our budget is 50 pruning epochs plus 25 fine-tuning epochs at batch 512. The
published CIFAR-10 results we quote for calibration use 160 epochs at batch 128
\citep{liu2021granet} or 250 epochs at batch 128 \citep{li2025east} --- close to an
order of magnitude more optimizer updates. Our absolute accuracies should
therefore be read as below what these architectures can reach, and no number in
this paper is a state-of-the-art claim. The internal comparison is unaffected,
because both arms are equally under-trained, but two second-order caveats survive:
we do not know whether the I.P.--BaCP gap would persist, grow or close under a
long schedule, and a short schedule generally flatters interventions that
accelerate early optimisation. Testing the gap at a full-length budget is the
single most valuable follow-up experiment.

\paragraph{Seed coverage.}
Three seeds were run on the headline cells only (Table~\ref{tab:seeds}); the rest
of the grid is single-seed. With a measured run-to-run spread of
\result{cifar10.noise.floor.lo}--\result{cifar10.noise.floor.hi} points, an
individual single-seed cell cannot on its own distinguish a small effect from
noise, and the objective ablation of Section~\ref{sec:results:ablation} is
single-seed throughout. What single-seed cells can support is a consistency
argument --- the sign of $\Delta$ agreeing across many independent cells is
evidence that no individual cell provides --- and that is how Table~\ref{tab:main}
should be read. It is not a substitute for replication
\citep{bouthillier2021variance}, and we do not present it as one.

\paragraph{Scope of the comparison.}
Finally, the paper compares BaCP against its own matched baseline, not against
other pruning methods. We quote published numbers for calibration
(Table~\ref{tab:published}), but we did not re-run RigL, GraNet, GraSP or SynFlow
under our protocol, and it would be wrong to read Table~\ref{tab:published} as a
ranking. Our arms are static; several of those methods are dynamic, and a
support-level intervention and an objective-level intervention are not
alternatives to each other. Whether they compose is an open question this paper
does not answer.
